\egasection{§13\quad 初等态射与 M. Artin 定理}
\source{\texttt{EGAVaiello.pdf}，PDF及印刷第49--53页}
\egas{13.1}{定义}预概形的态射$f\colon X\to Y$称为一个“初等态射”，如果$X$在$Y$上同构于形如$X'\setminus Z$的预概形，其中$X'\to Y$是光滑射影态射，具有几何连通的一维纤维，而$Z$是$X'$的闭子预概形，使得态射$Z\to Y$为满射、平展且次数恒定。

一个态射称为多重初等的，如果它是初等态射的复合。域$k$上的预概形$X$称为（在$k$上）多重初等的，如果其结构态射$X\to\operatorname{Spec}(k)$是多重初等的。\footnote{\textbf{英译编者注〔底本脚注63〕：}提醒Blass：询问Artin有关术语的问题。参见Milne第117页及SGA IV〔英译者注〕。}\cmt{D17-M046}{底本使用“elementary”和“polyelementary”，这里分别译作“初等”和“多重初等”。按第49页紧随其后的原作者说明，复合允许空族，这解释了恒等态射的情形；这里仍把该说明留作译注，没有冒充定义原段中已有的句子。脚注63未给出Milne书名或SGA IV的具体位置，故不补定引用。}
% END DOT-MIDDLE BATCH 09


% BEGIN DOT-MIDDLE BATCH 10: EGA第五卷_13.2归约_校订续译.tex
\egas{13.2}{定理（M. Artin）}设$k$是无限完美域，$X$是$k$上几何不可约的预概形，$x$是$X$的一个光滑点。则$x$具有一个由多重初等开邻域组成的基本邻域系。\cmt{D17-M047}{“多重初等”沿用定义13.1，指结构态射为初等态射的复合。底本定理本句未逐项重述有限型等环境假设，而下文立即使用仿射、拟射影与射影闭包的归约；本译稿不将这些略写解读为对任意无有限性条件的预概形都已给出证明。}

把$X$替换为$x$的一个给定邻域后，只须证明$x$在$X$中有一个多重初等开邻域〔底本英译者在句末注：最后这个字母是否正确？字迹不清〕。\cmt{D17-M048}{底本此句写“open elementary neighborhood”，而定理陈述是“open polyelementary neighborhoods”，下一段又按维数归纳并构造$U\to V$，故此处正文统一为“多重初等开邻域”，原用词留于此。初等结构态射的纤维为一维，不能直接把一般高维邻域称作在$k$上初等的。句末的字母难辨标记并未因此消除。}

对$X$的维数$n$作归纳，我们归约为证明：若$n>0$，则存在$x$的一个开邻域$U$及一个初等态射$f\colon U\to V$，其中$V$是$k$上的光滑概形（必然几何不可约，且维数为$n-1$）。（$n=0$的情形显然平凡，因为此时$X$同构于$\operatorname{Spec}(k)$；考虑到13.1并未排除空族态射的复合，它在$\operatorname{Spec}(k)$上是多重初等的。）我们应当在13.1中以某种方式提到这一点。\cmt{D17-M049}{此处保留原作者关于空复合的编写提醒。归纳中还须在$V$上取像点的适当邻域并拉回；底本没有在本段展开，本译文不把省略的归纳拼接冒充为原文已写出的步骤。}

首先要求$k$完美的必要性，在$n=1$时就已显现：此时取$X$的函数域$K$的典范射影正规模型作为$X'$（参见第II章第7段）。\footnote{\textbf{英译编者注〔底本脚注64〕：}EGA II，是的，7.4〔英译者注〕。}$k$完美保证$X'$在$k$上光滑（因为$X'$无论如何都是正则的），并且也保证带诱导约化结构的$Z=X'\setminus X$在$k$上平展。现在来处理一般情形，于是可以假设$n\geq2$。\cmt{D17-M050}{这一段关于曲线模型的写法隐含先在$x$附近作适当开限制，使所讨论的光滑曲线开集置于其正规射影模型中。并且定义13.1要求边界到基底满射，在域上即边界非空；若边界为空，还需在不含$x$处作适当删点。底本这里没有写出这些步骤，故指出略写而不声称已在正文补齐。}

显然可以假设$X$仿射，因而拟射影。再把$X$替换为一个射影闭包，就可以假设$X$射影，但始终须证明：每个邻域都包含一个允许初等态射$U\to V$的开邻域$U$。此外，把$X$替换为其正规化（它在$X$上有限，因而射影；引用待补），\footnote{\textbf{英译编者注〔底本脚注65〕：}我想是在EGA II〔英译者注〕。}并不会改变$x$的一个邻域，所以可以假设$X$正规；由于$k$完美，它因而在$k$上几何正规。\cmt{D17-M051}{“不改变一个邻域”按局部含义译出，并非断言正规化在整个$X$上是同构。原稿对正规化有限性的引用尚待定位，原脚注65的疑问语气保留。后续构造现已按底本接续收入，但这不等于该依赖链已经独立审定。}
% END DOT-MIDDLE BATCH 10


% BEGIN DOT-MIDDLE BATCH 11: EGA第五卷_13.3陈述_校订续译.tex
这一假设的好处是：$X$中在$k$上不光滑的点所成的集合$Z$，余维至少为2。选择一个射影浸入$i\colon X\to\mathbf P^r$。由于取各个$\mathcal O_{\mathbf P^r}(d)$（$d>0$）中在$x$处不消失的截面，其非零点开集构成$x$在$\mathbf P^r$中的一个基本邻域系，我们得到：$x$的每个邻域都包含一个形如$X\setminus Y$的邻域，其中$Y$是$X$中包含$Z$的闭子集，纯维数为$n-1$，且$x\notin Y$。\cmt{D17-M052}{底本把扭曲次数也写为$n>0$，与一直表示$\dim X$的$n$重名；这里改用$d$表示次数，未改变任何数值条件。“取截面”展开为取其非零点开集，以明确基本邻域系的对象。选择包含$Z$的边界是原稿在本段给出的结论，本译稿不添加一个新的选取证明。}

赋予$Y$诱导约化结构；由于$k$完美，$Y$的奇异点集的维数至多为$n-2$。扩大前面的集合$Z$，便得到一个维数至多为$n-2$的闭子集$Z\subset Y$，它包含$X$和$Y$的几何奇异点集。

证明的想法是：固定一个余维为$n$的线性子簇$C$，用$X$与$P$中包含$C$的余维$n-1$线性子簇$L$的交，来把$X$纤维化。为此，需要如下引理。

\egas{13.3}{引理}
沿用前述记号，设$X\supset Y\supset Z$是$\mathbf P^r_k$的闭子概形，维数分别为$n$、$n-1$及至多$n-2$，且$X\setminus Z$、$Y\setminus Z$光滑。$Z$的维数至多为$n-2$；必要时，\footnote{\textbf{英译编者注〔底本脚注66〕：}法文。}若$k$的特征为$p>0$，按第9节的办法把射影浸入$i\colon X\to\mathbf P^r$换成一个“倍数”$\phi_d\circ i$（$d\geq2$；底本在“任意”后又注“某个？”），则存在$\mathbf P^r$中一个余维为$n-1$的线性子簇$L_0$，具有以下良好性质：\cmt{D17-M053}{底本的包含关系印为$X\supset Y\subset Z$，与紧前一段$Z\subset Y$及维数限制不符，故订为$X\supset Y\supset Z$。底本重复的$Z$维数限制仍保留。复合符号在底本写作“$\phi_n^{i}$”，并以$n\geq2$为指标；依“替换射影浸入”及第9节的引用，正文明确记为$\phi_d\circ i$，以$d$区别于维数$n$。“any”与“[some?]”相冲突，正文保留不确定性，不断言任意次数均可。}
\begin{enumerate}
\item[a)] $L_0\cap Z=\varnothing$；
\item[b)] $L_0\cap X$光滑，维数为1；
\item[c)] $L_0\cap Y$光滑，维数为0。\cmt{D17-M054}{a）底本写“$\phi=\varnothing$”，这里直接采用空集符号，避免与前段的映射$\phi_d$混淆。}
\end{enumerate}


（注意：这里$k$是无限域，并不要求它完美。）现在假设这一引理成立，说明如何由此推出：存在$x$的一个开邻域$U$，包含在$X\setminus Y$中，并允许一个初等态射$U\to V$。\cmt{D17-M055}{本句启动的是利用引理的构造，其具体步骤见下段；引理的证明也在后文接续收入，但原稿的依赖与略写仍须分别核查。}
% END DOT-MIDDLE BATCH 11


% BEGIN DOT-MIDDLE BATCH 12: EGA第五卷_13.2初等邻域构造_校订续译.tex
存在$L_0$中的一个线性子簇$C$，它在$\mathbf P^r$中的余维为$n$，即在$L_0$中的余维为1，并且不与有限集$(L_0\cap Y)\cup\{x\}$相交。令$T=X\cap C$，则$T$是$X$中一个在$k$上平展的非空子概形，不含$x$，且与$Y$不交。\cmt{D17-M056}{底本两次写$L$，同段的$L_0\cap Y$及随后参数点所代表的子空间明确指向引理中的$L_0$，故统一为$L_0$。但仅仅避开所列有限集，并不足以从字面上推出$X\cap C$平展；还需适当横截选择$C$。此处保留原作者对所选$C$的性质的断言，并明确指出该选择步骤未在本段展开，不把任意满足避开条件的$C$都说成适用。}

另一方面，考虑$\operatorname{Grass}^{n-1}(P)$中对应于包含$C$的线性子簇的子概形$Q$。$Q$是一个维数为$n-1$的射影空间；特别地，它在$k$上光滑，维数为$n-1$。$L_0$对应于一个点$\xi_0\in Q(k)$。沿用此前引入的一般记号，考虑$X^{(n-1)}$沿浸入$Q\to\operatorname{Grass}^{n-1}(P)$的逆像$X_Q^{(n-1)}$，以及$Y_Q^{(n-1)}$、$T_Q^{(n-1)}$。后二者也是$X_Q^{(n-1)}$中互不相交的闭子概形。以$p,q,r$分别记这三个概形到$Q$的结构投影。\cmt{D17-M057}{底本先写参数点$\xi$，下页续文改写$\xi_0$，此处统一为$\xi_0$。列举第三个逆像时底本只印出带上下标的$Q$而漏掉对象字母；紧随其后的$T_Q^{n-1}\simeq T\times_k Q$及边界构造确定它是$T$的逆像。括号中的上标$(n-1)$表示余维，排版整理不改变原含义。}

由假设，$p$在$\xi_0$上方的点处光滑，$q$在$\xi_0$上方的点处平展；$r$也是如此，因为$T_Q^{(n-1)}$无非就是$T\times_k Q$（在$Q$上同构）。最后，态射$p$固有；又由于$X$几何连通，$p$的纤维几何连通（Bertini定理）。\cmt{D17-M058}{这一连通性结论按原稿引用保留。底本并未在此写明所用Bertini定理的完整假设或引用位置，本译稿不以“$X$几何连通”单独替代所需的线性截面连通性论证。}

因此，存在$\xi_0$在$Q$中的一个开邻域$V$，使得
\[
X'=X_Q^{(n-1)}\big|_V
\]
在$V$上固有、光滑，并有几何连通纤维。由于$\xi_0$处的纤维正是$X\cap L_0$，其维数为1，可以假设$X'$在$V$上的所有纤维都为一维。最后，将$V$取得充分小，可以假设$Y_Q^{(n-1)}|_V$及$T_Q^{(n-1)}|_V$在$V$上平展，因而二者的和概形$Z'$（它可识别为$X'$的一个闭子概形）在$V$上平展。于是，令
\[
U=X'\setminus Z',
\]
便得到一个初等态射$U\to V$。\cmt{D17-M059}{底本原句为“an open neighborhood $U$ of $Q$ in $X$”，但紧接着把族限制到$V$，上下文始终是在参数点附近缩小基底，故正文订为“$\xi_0$在$Q$中的开邻域$V$”；原“fiber of $Q$”同理订为$\xi_0$处的纤维。定义13.1所需射影性可由原射影关联族限制得到；边界的满射性和恒定次数还须使用非空$T$及必要的缩小。本段原文没有把这些核对全部写出。}

但$U$也在
\[
X''=X_Q^{(n-1)}\setminus
\bigl(Y_Q^{(n-1)}\cup T_Q^{(n-1)}\bigr)
\]
中开；$X''$是$X\setminus(Y\cup T)$在$X_Q^{(n-1)}$中的逆像。另一方面，$X''\to X\setminus(Y\cup T)$是同构，因为
\[
X_Q^{(n-1)}\setminus T_Q^{(n-1)}\longrightarrow X\setminus T
\]
是同构。因此，$U$可识别为$X\setminus(Y\cup T)$的一个开子集；它包含$(L_0\cap X)\setminus(Y\cup T)$，因而包含$x$。这就是所需的、包含在$X\setminus Y$中的$x$的邻域。\cmt{D17-M060}{底本在$X''$的公式中写$Z_Q^{n-1}$，但随即说明它是$X-Y-T$的逆像，且本段移去的第二个边界一直是$T$，故订为$T_Q^{(n-1)}$。底本又写$U$包含整个$L_0\cap X$；因$U$已删去边界，这应为删去$Y\cup T$后的部分，正文据构造修正。最后“因而包含$x$”还使用$x\in L_0$：底本引理13.3的陈述漏明这一点，下一段证明却明确取经过$x$的线性子簇。此项前后依赖在这里指出，不能只凭前述引理字面推出。}
% END DOT-MIDDLE BATCH 12


% BEGIN DOT-MIDDLE BATCH 13: EGA第五卷_13.3证明归约_校订续译.tex
现在只剩下证明引理13.3。按通常的做法，只需证明经过$x$、余维为$n-1$的一般线性子簇（记为$L$）具有性质a)、b)、c)。\cmt{D17-M061}{此处“一般”对应原文generic；后文用相应参数空间的一般点说明其含义，不是任意线性子簇。经过$x$这一条件在本段明确出现，与前文指出的引理陈述遗漏相呼应。}

性质a)，以及b)、c)中的维数部分，可直接由2.3（已经在第12节重新审视并作修正）得到：如第8节中已用过的一次论证那样，将它应用于参数化经过$x$的直线的射影空间，以及$Z$经以$x$为中心的锥投影在该射影空间中的像。\cmt{D17-M062}{底本此处印“passing through $X$”，但同句投影中心及整段指定点均为$x$，故订为$x$。“[dans ledit]”是底本保留的法语，指前述空间；正文译出这一指代。原文称维数结论直接来自2.3，本译稿不另补写其未展开的推导。}

［此外，将用这种方法所得的某些结果明确写出，或许也有益处：这些结果涉及由满足经过一个固定线性子簇这一条件的线性子簇所截出的线性截面。可放在正文中，也可另设一节。］\cmt{D17-M063}{方括号内是原稿的编写建议，不是本译稿新增的证明或已完成的补充章节。}

对于b)、c)中的光滑性，由于a)，可以分别用$X\setminus Z$和$Y\setminus Z$代替$X$和$Y$；它们都是光滑的。于是问题归结为证明以下断言。

设$f:X\to P$是一个无分歧态射，$X$在$k$上光滑，且$x\in P(k)$不属于$X$中任何维数小于$m$的分量的像（不可约分量？）。\cmt{D17-M064}{括号中的疑问“irreducible component?”来自底本，照录而不替作者作定论。底本后半句印“let $X\in P(k)$”；因$X$已是态射的定义域，此处应为点$x$，正文据前后文整理，且将该点的说明前移以免悬空。}
令$\eta$为$\operatorname{Grass}^{m}P$中由经过$x$的余维$m$线性子簇$L$组成的子Grassmann簇的一般点。则相应的一般线性截面$X^{(m)}_\eta$在$\kappa(\eta)$上光滑：至少在$k$的特征为零时如此；在另一种情形中，则须将$f$换成$Q_n f$，其中$n$为不小于2的整数。\cmt{D17-M065}{底本写“$X^{(m)}$ is smooth over $k$”，未在这里印一般点下标；底本后续第52页明确写$X^{(m)}_\eta$。正文用“一般线性截面”表达其上下文所指，而不把整个关联族都断言为光滑。底本原写“smooth over $k$”。这里的一般纤维定义在$\kappa(\eta)$上，正文据此订正光滑性的基域；不能把沿参数空间一般点的基变换偷换为对原域$k$的光滑态射。光滑态射要求局部有限表示，而例如$\Spec k(t)\to\Spec k$并不满足这一要求。此项基域校正不补齐所引Bertini结论的全部假设或缺文，详见核查项D17-M-E01。底本此处确实印$Q_n f$，紧接的下一段改用$\Phi_n f$；此项符号不一致保留待核，不静默将二者认作同一映射。这里的$n$又用作该替换的整数指标，与段首维数记号重名，照录并提示。}
% END DOT-MIDDLE BATCH 13


% BEGIN DOT-MIDDLE BATCH 14: EGA第五卷_13.3一般截面归纳_校订续译.tex
这是对第9节的一项补记，而它本身又来自第8节后面的补记。\cmt{D17-M066}{底本在这些位置使用regret、regrets，按上下文是希望补入先前章节的结果，故译为“补记”，并非已正式编号的勘误。}
沿用8.8的记号，并假设$X$不可约；对于我们正在讨论的$p^0$问题（？），这样做是允许的。\cmt{D17-M067}{“$p^0$-problem (?)”为底本原有的难辨记号与疑问，未擅自还原成某一性质。本译稿所用底本在8.7后没有完整的8.8条目，故此引用并不表示8.8缺文已经恢复。}
如果
\[
\operatorname{codim}T\geq2,
\]
或者$h^{\mathrm{sing}}\to T$一般平展，那么对于经过$x$的一般超平面$H_\eta$，截面$X^{(1)}_\eta$是光滑的，维数为$n-1$；这里须除去$f(X)=\{x\}$的情形（于是$n=0$）。\cmt{D17-M068}{底本此处印$h^{\mathrm{sing}}\to T$，第52页后续证明改印$X^{\mathrm{sing}}\to T$，本译稿不将两种记号静默统一。原句“the hyperplane $H_\eta$ passing through a generic $x$”词序含混；上一段及下一页均把$x$固定、令$\eta$为经过$x$的超平面参数空间的一般点，正文据此整理为“经过$x$的一般超平面”，不另要求$x$是$X$或$P$的一般点。末句原印$f(x)=\{x\}$，按其所说$X$零维的例外及下一页对所有$z\in X$的讨论，应为$f(X)=\{x\}$，故正文订正。维数$n$与下句替换指标$n$重名，仍照底本保留。}
一般平展这一条件在$k$的特征为零时自动满足；或者，在将$f$换成$\Phi_n f$、其中$n\geq2$之后满足，参见第9节。\cmt{D17-M069}{底本把这一条件说明插在前句中，正文为避免长括号而移作紧随其后的句子。这里保留其断言与引用，不把第9节的适用假设或所引结论视为已独立验证；前文原文的$Q_n$与此处$\Phi_n$之间的记号不一致仍待核。}
暂且承认这一结果，它就解决了上述补记中$m=1$的特殊情形。对$m$作归纳即可得到一般$m$的情形：注意，允许作基变换后，$X^{(m)}_\eta$可以先取一个经过$x$、余维为$m-1$的$L'$，再取由$H$给出的线性截面而得到；$L'$和$H$对于这些条件是相互独立的一般选择。\cmt{D17-M070}{底本原句在“taking an $L'$”之后把后续动作放进长括号，语法不完整；这里仅按它明确列出的两步整理语序。所取的是相应关联参数空间上的一般点，而非任意两个线性子空间。没有由此声称一般纤维无需任何基变换即可相等。}
换成通常的严谨说法，就是把自己置于参数化二元组$(L',H)$的概形的一般点处，并取由$H$给出的线性截面；归纳假设提供了所需的光滑性。\cmt{D17-M071}{底本括注末尾为“by taking the linear sections by $H$ which is smooth by inductive assumption”，其中which的指代不清。归纳结构中，应先用归纳假设处理余维$m-1$的截面，再用刚承认的$m=1$情形处理最后一个超平面截取，不能仅凭归纳假设就断言最后的截面光滑。正文保留其归纳说明，脚注明确这一步仍依赖尚待证明的$m=1$结果。}

这种类型的论证此前已用来推广2.6，例如推广到任意余维$m$的线性截面。值得把它在一般情形下一次性清楚写出，以便以后引用，而不必每次都进入一份正式表述中那些有些繁重的细节［法语；此处有空缺］。\cmt{D17-M072}{底本末尾为“a bit heavy [Fr] [\quad] of a presentation (enforme)”。保留其未辨空缺；enforme按“en forme”的语义译为正式表述。这一段是作者关于编写方式的说明，不意味着本译稿已经补成了一项一般定理及其证明。}
% END DOT-MIDDLE BATCH 14


% BEGIN DOT-MIDDLE BATCH 15: EGA第五卷_13.3超平面情形证明_校订续译.tex
还须证明上面宣布的8.8的推论，即$m=1$的情形。如果$\operatorname{codim}T\geq2$，另一方面，参数化满足$x\in H_\xi$的超平面的子空间$Q$是对偶射影空间$P^\vee$中的一个超平面，余维为1；因此$Q$的一般点$\eta$不可能属于$T$，于是便得到了所需的结论。\cmt{D17-M073}{底本在结句保留法语“on gagne [Fr]”，随后附有英译疑问“(we are done?)”，此处译出其完成论证之意，并保存该疑问。这里用到的是$T$作为闭的坏参数集所具有的含义：若$Q$的一般点属于$T$，则整个$Q$包含于$T$，与余维大小冲突；本译稿没有恢复8.8缺失的完整假设。}

在$\operatorname{codim}T=1$的情形，由于$T$不可约，只有当$Q=T$时，才可能有$\eta\in T$。\cmt{D17-M074}{底本原句为“we cannot have $\eta\in T$ under $Q=T$”，按字面与此处论证相反。由$\eta\in T$得$Q\subseteq T$，而二者为同维不可约闭子簇，所以必须相等；后文也明确在$T=Q$的情形下继续论证。正文据此订为“只有当……时才可能”，原连接词保存在本注。}
用几何语言来说（可以假设$k$代数闭），这意味着：对每个$z\in X$，$X$在$z$处的切空间，确切地说，是该切空间经$f'_z$所得的像，经过$x$。\cmt{D17-M075}{底本的代数闭括注没有正常闭合；这里整理括号，不补造缺失论证。原“or also plutot [Fr]”按法语plutôt译为“确切地说”。切空间在此按原射影对偶论证理解，不将“经过$x$”当成抽象向量空间的无坐标陈述。原段没有逐项重述对偶簇$T$及切空间像的定义，适用条件仍须与前文核对。}
我们来证明，在$X^{\mathrm{sing}}\to T$一般平展、也就是处于所引用的8.8条件之下时，这种情形不能发生，除非$f(X)=\{x\}$，此时$X$为零维。\cmt{D17-M076}{底本两次使用unless：“cannot happen unless $X^{\mathrm{sing}}\to T$ is generically étale, i.e. unless we are under the conditions of 8.8”。但紧前原段把一般平展列为获得结论的充分条件，紧随其后的论证也正是在该条件下引用8.7(c)得出矛盾；正文依此订为“在……条件下不能发生”。这一校订只处理连接词与上下文的逻辑方向，不补齐8.8的缺失条件，也不宣称该对偶结论已获独立证明。此处原印$X^{\mathrm{sing}}$，上一页印$h^{\mathrm{sing}}$，记号不一致照录待核；不能未经核对就将它解释为当前光滑$X$的奇异点集。底本$f(x)=\{x\}$据同段对所有$z$及$X$维数的讨论订为$f(X)=\{x\}$。在原先$f$无分歧的有限型背景中，像为一点对应零维情形；并非声称任意常值态射的定义域都零维。}

事实上，由8.7(c)（它本质上表达了$X$与其“对偶”$T$之间关系的对称性），对于几乎所有$z\in X(k)$，$f(z)$与$T$在某一点处的切空间正交。由于$T=Q$，这就意味着$f(z)$与$Q$正交，故$f(z)=x$，从而$f(X)=\{x\}$。这证明了我们的补记，因而也证明了13.2。\cmt{D17-M077}{底本后半句语序破碎：“at a certain point this (since $T=Q$) orthogonal to $Q$, where $f(z)=x$”；正文据其显式等式整理为连续推论。在射影对偶意义下，$Q$是所有经过$x$的超平面的参数超平面，其正交点就是$x$。末句保留的是原作者关于论证完成的表述，不是译者对全部依赖链的独立审定。“几乎所有”在此是代数几何中一般点意义的说法，不是概率或测度断言；从一般$z$的结论推广到整个像所需的稠密性及闭性使用在原稿中未展开。尤其8.7(c)和8.8的原文残缺仍未由本段恢复。}
% END DOT-MIDDLE BATCH 15


% BEGIN DOT-MIDDLE BATCH 16: EGA第五卷_13.4有限域与连通性_校订续译.tex
\noindent\textbf{附记。}
如果把$x$换成一个维数大于0的线性子簇$C$，并要求$H$经过$C$，上述论证就不能照样进行。\cmt{D17-M078}{底本第一句写“a linear sub-variety $X$”，但同句后半和随后讨论固定线性子簇均用$C$，这里据上下文统一为$C$；不改变原先源概形$X$的含义。}

事实上，没有理由认为［本句连接关系待核］（例如，取$\dim C=r-2$，这是$H$仍能实际变化时所允许的最大维数），“不作假设（？）”，$T$包含直线$C'$；例如，在任意特征下取$P$中的一个非奇异二次超曲面。\cmt{D17-M079}{底本可辨原句为“indeed, there is no reason to suppose (taking for example dim (= r − 2 the greatest possible for which H can still effectively vary) without supposing (?) that T contains the straight linear $C'$ taking for example a non-singular quadric in p in any characteristic.” 原有“dim (”处按同段对象整理为$\dim C$；但“no reason to suppose”与“without supposing (?)”之间关系不明，不能据此确定所否定的包含关系。正文逐项保留全部可辨内容与原问号，不把该残句改写为已经证明的反例。$C'$应如何精确定义，底本本段也没有明确写出；不擅自增补其定义或具体二次超曲面。}

不过，对于$\varphi_n f$，其中$n\geq2$，这样的现象或许不再发生；我们可以把这个问题作为第9节的一条注提出。\cmt{D17-M080}{底本本处使用小写$\varphi_n$，此前用$\Phi_n$或$Q_n$，记号差异保留待核。“或许”与“可以提出问题”均为原作者的未定语气，不译成肯定定理。}
\egas{13.4}{注}

\noindent a) 我们已经指出，对于13.2的成立，$k$是完美域这一假设是必要的。相反，$k$无限这一假设很可能是多余的、过强的。我们不打算针对$k$有限的情形给出专门的论证，只指出，在这种情形下，将13.2应用于$k$的代数闭包，再用通常的论证，可以找到$k$的一个有限扩张$k'$，使得$X_{k'}$中位于$x$上方的点具有相对于$k'$的多重初等开邻域。\cmt{D17-M081}{底本印“the point of $S_{k'}$ over $x$”，而13.2的概形为$X$，本条没有定义$S$，故正文订为$X_{k'}$。基变换后$x$上方不一定只有一个点，中文不额外引入唯一性断言；该句的逐点范围及有限扩张选择的下降细节，原稿没有展开。本条只是作者关于有限域情形的讨论，不以其推测断言有限域上的原结论已经成立。}

\noindent b) 如果在13.2中舍弃$X$几何不可约这一假设，则结论不能保持原样，因为一个多重初等代数概形是几何不可约的。不过，结论仍以一种较弱的形式成立：在定义13.1中删去“连通”一词便得到这种形式；我们所给的证明表明了这一点。\cmt{D17-M082}{定义13.1原文要求紧化族具有“几何连通的一维纤维”。本条“删去连通”沿用作者字面指示，是对初等态射定义的修改，不能仍称作完全相同的原定义。正文保留原作者对证明的判断；本译稿没有逐项重新验证去掉几何不可约后所有构造和边界条件。零维情形、空复合及非空条件须连同13.1、13.2前文理解，不外推为任意概形的断言。}
% END DOT-MIDDLE BATCH 16


% BEGIN DOT-MIDDLE BATCH 17: EGA第五卷_13.4有限点集邻域_校订续译.tex

\noindent c) （或可收入13.2的陈述。）沿用13.2的记号，设$\Phi$是$X$的一个有限子集，由$X$的光滑点组成，并假设$\Phi$包含在$X$的一个仿射开集中。则$\Phi$具有由多重初等开邻域组成的基本邻域系。显然，可以假设$\Phi$由闭点组成。

证明本质上相同，只须将13.3的表述略加改变如下：存在$P^r$中一个余维为$n$的线性子簇$C$，它不与$\Phi\cup Y$相交，并且对于每个$x_i\in\Phi(k)$，由$C$与$x_i$张成的余维为$n-1$的线性子簇$L_i$具有13.3的性质a)、b)、c)。\cmt{D17-M083}{底本写“a linear subvariety $C$ of $P^r$ of codimension not meeting $\Phi\cup Y$”，在“codimension”之后漏去数值。同句规定$L_i=\langle C,x_i\rangle$的余维为$n-1$，且$x_i\notin C$，故$\dim L_i=\dim C+1$，所缺数值应为$n$；这也与第50页引理13.3之前的构造相符。底本“the linear subvariety of $L_i$”的多余“of”在正文删去。}\cmt{D17-M084}{底本先说可假设$\Phi$由闭点组成，随后却写$x_i\in\Phi(k)$。闭点未必是$k$-有理点；此处照存记号，不擅将闭点改成有理点，也不增写未给出的基域扩张与下降论证。该过渡仍待核。}

为验证这一点，注意只须验证一般的$C$具有上述性质，因为对于这样的一个［原文缺词：???????］，每个$L_i$在经过$x_i$、余维为$n-1$的线性子簇$L$之中都是一般的，\cmt{D17-M085}{此处底本写“among the $L$ of codimension $n$ passing through $x_i$”。然而同段的$L_i$及13.3中的$L_0$的余维均为$n-1$，故将这里的$n$订为$n-1$。对于固定的有理点$x_i$，每个经过$x_i$的这种$L$都可由其中不经过$x_i$的一个超平面$C$与$x_i$张成；这是与该余维数相容的线性代数说明，不是对有限点集情形中全部光滑性条件的独立证明。七个问号在底本中实际存在，其缺词不予恢复。}因此可以应用13.3的最初形式（或者至少应用我们已经证明的形式，即［法文标记］：每个经过$x$且充分一般的$L_0$都具有性质a)、b)、c)）。\cmt{D17-M086}{底本末句保留“which was [Fr]”的编写痕迹；正文保留其法文标记，不臆造缺失的法语文字。这里“每个……充分一般的”是原文已证明形式的措辞，不应读成每一个经过$x$的线性子簇均满足这些性质。}
% END DOT-MIDDLE BATCH 17


% BEGIN DOT-MIDDLE BATCH 18: EGA第五卷_13.4一般基底_校订续译.tex

\noindent d) 按照7.1所说明的方法，当把基域$k$换成一般的基预概形$S$时，可以给出13.2的若干变式。我们指出下面一点（不作证明）：\cmt{D17-M087}{底本开头写“a general base $Y$ prescheme”，随后写$f:X\to Y$，但紧接着所有基底条件均使用$S$，而末句又以$Y$表示$X$中的闭子概形。为区别这两个对象，正文统一将前两处作为基底的$Y$改为$S$，保留末句的闭子概形$Y$。这只是消除记号冲突，不是修改基底假设。}

设$f:X\to S$是一个平坦射影态射，其纤维几何不可约且满足$(R_2)$；设$S'$是$X$的一个在$S$上有限的子概形，$s\in S$。假设对于每个位于$s$上方的$x\in S'$，纤维$X_s$在$x$处在$k(s)$上光滑。\cmt{D17-M088}{底本写“$x\in S$”，但后文取$s$的邻域，并用$x$表示位于$s$上方的$S'$中的点，因此这里改为$s\in S$。纤维条件$(R_2)$照底本保留；原文没有在本项说明它是否意指几何意义下的条件，不能擅改为$(R_1)$或更强条件。}
则存在$s$的一个开邻域$U$，以及$X|_U$中$S'|_U$的一个开邻域$V$，使得
\[
V\longrightarrow U
\]
是多重初等态射。\cmt{D17-M089}{此为作者明说“without proof”的断言。所引7.1在有限剩余域的情形明确允许先作平展满射基变换；本项却只写原基底中的开邻域$U$，没有说明何时可以免去该基变换，也没有补列剩余域条件。本译文不以7.1的引用消除这一问题，亦未独立验证本项所述的一般性。}

若$Y$是$X$的一个不与$S'$相交的闭子概形，而$Z$为$Y$在$S$上不光滑的点组成的集合，且［原句连接语残缺，待核］
\[
\dim Z_s\leq\dim X_s-2,
\]
则在上述结论中，还可以取$V\subset X-Y$。\cmt{D17-M090}{底本末句可辨部分为“such that the set $X$ of points where $Y$ is not smooth over $S$ $\dim Z_s\leq\dim X_s-2$ then we may above tale V containe din $X-Y$”。这里点集先写$X$，不等式却用未另定义的$Z_s$；按该句的唯一非光滑点集及其纤维记号，正文将点集名订为$Z$，以免与全空间$X$混淆。“tale V containe din”按上下文读作“take $V$ contained in”。底本在“over $S$”之后确实缺少语法连接；正文保留残缺标记，不额外补出对所有基底点的量词、闭性假设或证明。}
% END DOT-MIDDLE BATCH 18


% BEGIN DOT-MIDDLE BATCH 19: EGA第五卷_13.4拓扑结构_校订续译.tex

\noindent e) 13.2之所以值得关注，原因之一在于多重初等代数概形$U$的拓扑结构特别简单。\cmt{D17-M091}{底本写“the elementary algebraic schemes $U$”。13.2的邻域为多重初等概形，13.1也以结构态射为初等态射之复合来定义多重初等概形；本项随后讨论一般维数$n$以及自由群的逐次扩张。因此此处正文用“多重初等”，不把$U$误限为结构态射本身是初等态射的相对一维情形。}
例如，若基域为复数域，并以$U^{\mathrm{en}}$表示与$U$关联的解析空间，则同伦群
\[
\pi_i(U^{\mathrm{en}})=0\qquad(i\geq2),
\]
而$\pi_1$是自由群的逐次扩张。\cmt{D17-M092}{底本先写“if $Y^{\mathrm{en}}$ denotes the analytic space associated to $U$”，后文却始终使用$U^{\mathrm{en}}$，故将前一个$Y$订为$U$，上标“en”照存。底本同伦群指标写“$i\ne1$”；按同伦群的正次数范围，这就是$i\geq2$。$\pi_0$通常是连通分支集合而非这里的同伦群；正文不将其当作群来写等于零。未增写或指定原文未给出的扩张列。}

因此，$U^{\mathrm{en}}$是一个为$\pi_1$分类的“$K(\pi_1,1)$空间”；更确切地说，它的普遍覆盖空间同胚于$\mathbb C^n$，从而当然是可缩的，而这一覆盖是一个群为$\pi_1$的“普遍主纤维化”。\cmt{D17-M093}{底本写“homeomorphic to $C^n$”，这里$C$按复数域排为$\mathbb C$，$n$沿用本节维数记号。“同胚”仅是拓扑意义，原文没有声称与$\mathbb C^n$双全纯。普遍覆盖可缩比“同胚于$\mathbb C^n$”弱，不能由前者反推后者；本项未附这项较强断言的证明。本译文保存原断言，不宣称已经独立验证其拓扑论证。末句的覆盖及群作用是作者所谓“普遍主纤维化”的语境，不另行推定一个未给出的代数主丛构造。}
% END DOT-MIDDLE BATCH 19

